\subsection{Многотактная схема вычисления функции}

\subsubsection{Постановка задачи}

Для вычисления значения функции \eqref{eq:func} необходимо произвести три арифметические операции:
\begin{itemize}
    \item вычисление кубического корня числа $a$;
    \item вычисление квадратного корня числа $b$;
    \item суммирование результатов двух предыдущих операций.
\end{itemize}

Все перечисленные операции были реализованы в разделе \ref{sec:multi_cycle_modules}.

Дополним граф зависимости модулей, чтобы реализовать вычисление функции. Результат представлен на рисунке \ref{fig:module_dependencies_function}.

\begin{figure}[h]
    \centering
    \begin{tikzpicture}[
        fsm/.style={
            draw, rectangle, thick, fill=blue!10, 
            minimum width=4cm, minimum height=1cm, 
            align=center, rounded corners
        },
        module/.style={
            draw, rectangle, thick, fill=green!10, 
            minimum width=3.5cm, minimum height=1.2cm, 
            align=center
        },
        resource/.style={
            draw, rectangle, thick, fill=orange!15, 
            minimum width=3.5cm, minimum height=1cm, 
            align=center
        },
        data/.style={
            -{Stealth[scale=1.2]}, thick, draw=black
        },
        node distance=1.5cm and 1.5cm
    ]

    % --- Узлы ---
    
    % Управляющий автомат
    \node[fsm] (fsm) {Модуль\\ результата функции};

    % Модули алгоритмов
    \node[module, below=1cm of fsm, xshift=-2.5cm] (cbrt) {Модуль\\кубического корня};
    \node[module, below=1cm of fsm, xshift=2.5cm] (sqrt) {Модуль\\квадратного корня};

    % Общие аппаратные ресурсы
    \node[resource, below=3cm of fsm] (adder) {Модуль сумматора};
    \node[resource, below=3cm of fsm, xshift=-5cm] (mult)  {Модуль умножителя};

    % --- Связи ---

    \draw[data] (fsm) -- (cbrt);
    \draw[data] (fsm) -- (sqrt);
    \draw[data] (fsm) -- (adder);

    \draw[data] (cbrt) -- (adder);
    \draw[data] (sqrt) -- (adder);
    \draw[data] (cbrt) -- (mult);
    \end{tikzpicture}
    \caption{Граф зависимости модулей для вычисления функции.}
    \label{fig:module_dependencies_function}
\end{figure}

\subsubsection{Модуль результата функции}

Граф переходов для модуля результата функции представлен на рисунке \ref{fig:func_result_fsm}.

\begin{figure}[h]
    \caption{Граф переходов для модуля результата функции.}
    \label{fig:func_result_fsm}
    \centering
\end{figure}

Код модуля результата функции представлен в листинге \ref{lst:func_result} приложения \ref{app:listings}.

\subsubsection{Область допустимых значений функции}

\dots